% Generated draft: requires mathematical and source review.
\documentclass[11pt]{article}
\usepackage[margin=25mm]{geometry}
\usepackage{fontspec}
\setmainfont{Latin Modern Roman}
\usepackage{amsmath,amssymb,mathtools,unicode-math}
\setmathfont{Latin Modern Math}
\usepackage{enumitem,needspace,xurl,hyperref}
\setlength{\parindent}{0pt}
\setlength{\parskip}{5pt}
\setlength{\emergencystretch}{3em}
\newcommand{\N}{\mathbb{N}}
\newcommand{\Z}{\mathbb{Z}}
\newcommand{\Q}{\mathbb{Q}}
\newcommand{\R}{\mathbb{R}}
\newcommand{\C}{\mathbb{C}}
\newcommand{\F}{\mathbb{F}}
\newcommand{\A}{\mathbb{A}}
\newcommand{\PP}{\mathbb{P}}
\newcommand{\E}{\mathbb{E}}
\newcommand{\eps}{\varepsilon}
\newcommand{\dd}{\,\mathrm{d}}
\begin{document}
\section*{UnsolvedMath ID: 11100060}
\textbf{Source problem code:} catalogue entry
\subsection{Definitions and mathematical statement}
At a chosen prime \(p\), let \(BP\) denote Brown--Peterson theory, conventionally with
\[
 BP_*\cong\Z_{(p)}[v_1,v_2,\ldots],\qquad |v_j|=2(p^j-1),
\]
for a specified standard generator convention. For \(x\in BP_n(X)\), \(v_n\)-torsion means \(v_n^r x=0\) for some integer \(r\geq1\), with multiplication taking place in the corresponding higher homological degree.

\emph{The Johnson question (source formulation).} If \(X\) is a space and \(x\in BP_n(X)\), is \(x\) not \(v_n\)-torsion? The same \(n\) occurs as the homological degree and the generator index; it is not replaced by an arbitrary independent height. The literal phrase ``\(x\) is not torsion'' includes the zero element, which is killed by every power. Its usual meaningful no-nonzero-torsion interpretation would be
\[
 \forall X\ \forall x\in BP_n(X)\ \forall r\geq1,
 \qquad v_n^r x=0\ \Longrightarrow\ x=0.
\]
This interpretation is made explicit rather than silently excluding zero. The source does not state the prime quantifier, generator convention, allowed \(n\), or whether degree zero uses \(v_0=p\); the displayed coefficient notation uses the usual positive-height range \(n\geq1\) conditionally. No finite-complex, connectivity, or localization hypothesis on \(X\) is added beyond the intended \(BP\) setting.

Hovey proposes the question as a test for a powerful \(BP\)-analogue of Lannes's theory of unstable algebras over the Steenrod algebra. He credits Boardman--Johnson--Wilson with constructing such a framework and says it had not resolved this question in his context. The summary distinguishes this literal \(BP\)-homology-of-spaces question from similarly named questions about \(v_n\)-torsion in homotopy of Moore or finite complexes. Those different groups are not substituted here, nor is homology replaced by \(BP\)-cohomology.
\subsection{Short English statement}
Ask whether multiplication by every power of \(v_n\) kills no nonzero element in degree \(n\) Brown--Peterson homology of a space. The source ties the degree and height indices together but omits prime/range conventions and literally includes zero; its unstable-theory motivation is retained.
\subsection{Sources}
\begin{itemize}
\item Local UnsolvedMath registry, permanent ID 11100060, source code catalogue entry; proposer Mark Hovey, ``Unstable homotopy theory,'' item 1, explicitly titled the Johnson question.
\item \emph{Hovey's Algebraic Topology Problem List}, surviving mirror: \url{https://www-users.cse.umn.edu/~tlawson/hovey/}. Recorded access 2026-07-29; extraction html-ordered-list. The mirror has 65 top-level items across nine category pages.
\item The statement names Lannes and Boardman--Johnson--Wilson; the summary also mentions Johnson--Wilson and unstable \(BP\) modules. No full bibliography or precise theorem beyond the quoted question is supplied.
\item Supplied status open and OPEN-TRIAGE. The summary reports no verified resolution found and cautions about similarly named Moore-complex/homotopy questions. Those are supplied search assessments. Current primary-literature status and redistribution terms remain NEEDS\_REVIEW. The registry assigns L4 whereas background records default L3.
\end{itemize}
\end{document}
